\chapter{纯 Python 零依赖 Qwen 模型真实权重推理}

\section{Lab 07b：亲手实现 Qwen2.5 的推理，并与官方逐位对拍}

\begin{noteBox}[核心导读与背景]
这是整条学习路线里\textbf{"从用库的人变成懂原理的人"}的关键一步。

对应理论：第 4、5、6、7、8 篇全部。
前置：完成 Lab 05（你已经在自己训练的模型上写过同样的结构）。
依赖：PyTorch、transformers、safetensors、huggingface\_hub（见根目录 \texttt{requirements.txt}）。第一次运行会下载约 1GB 权重，国内网络可先 \texttt{export HF\_ENDPOINT=https://hf-mirror.com}。
\end{noteBox}

\section{规则}

\texttt{transformers} 只允许用来做两件事：\textbf{分词}，以及\textbf{作为标准答案和我们对拍}。
模型的前向传播、RoPE、GQA、KV Cache、重复惩罚全部自己写，直接读取 \texttt{model.safetensors} 里的原始张量。核心的 \texttt{QwenFromScratch.forward}（含 RMSNorm / RoPE / GQA 注意力 / SwiGLU / KV Cache）约 80 行。

\section{运行}

\begin{lstlisting}[language=bash]
.venv/bin/python lab07b_qwen_from_scratch/qwen_from_scratch.py
\end{lstlisting}

\section{本机（RTX 5060 Ti）一次运行的关键结果}

实验 1\textasciitilde{}5 的数值是确定性的，并且由脚本断言；实验 6 是计时。
\begin{noteBox}[核心导读与背景]
{}[计时] 实验 6 的计时数字是本机某一次运行的\textbf{量级示例}，前提是显卡空闲、已预热；GPU 降频或有其它负载时可能慢 40\% 以上。请看\textbf{趋势和比例}，以你自己运行的输出为准。
\end{noteBox}

\begin{center}
\small
\begin{tabularx}{\textwidth}{p{3.5cm} X}
\toprule
\textbf{实验} & \textbf{结果} \\
\midrule
1. 读 config.json & GQA 分组 7（14 个 Q 头 / 2 个 KV 头），每 token KV 仅 \textbf{12 KB} \\
2. 参数分布 & 词嵌入 \textbf{28\%}、FFN 64\%、注意力只有 9\%——小模型的词表开销惊人，所以共享 LM Head \\
3. BPE 分词 & "大模型推理的瓶颈是显存带宽。" $\rightarrow$ 11 个 token：\texttt{大/模型/推理/的/瓶颈/是/显/存/带/宽/。} \\
4. logits 对拍（FP32） & 最大误差 \textbf{1.5e-5}，每个位置 top-1 全部一致 \\
5. 生成对拍 & 纯贪婪逐 token 一致；并且发现 HF 的 \texttt{generate(do\_sample=False)} \textbf{偷偷使用了} \texttt{generation\_config.json} 里的 \texttt{repetition\_penalty=1.1}，我们按第 7 篇公式手写惩罚后再次逐 token 一致 \\
6. 性能（BF16） & Prefill 1024 token $\approx$ 80 ms（约 1.25 万 token/s）；Decode 约 20 ms/token，只达到带宽下界 2.6 ms 的\textbf{一成左右}（本机多次运行在 7\%\textasciitilde{}14\% 之间波动，取决于 GPU 频率与负载） \\
6. 朴素 vs Cache & 上下文越长差距越大：128 时两者接近（约 0.7\textasciitilde{}1.7x，此时单步耗时被 kernel 启动等固定开销主导，谁快取决于当时的噪声），2048 时约 \textbf{9\textasciitilde{}11x} \\
\bottomrule
\end{tabularx}
\end{center}

\subsection{为什么 Decode 只达到理论上限的一成左右？}

权重只有 1GB，按带宽算每 token 只要 2.6ms，但我们的实现每一层要启动十几个小 kernel（24 层就是几百个），每个 kernel 启动都有微秒级开销，逐元素算子还在反复读写显存，GQA 的 \texttt{repeat\_interleave} 甚至把 KV 复制了 7 份。
\textbf{这九成左右的差距，就是 vLLM / llama.cpp / TensorRT-LLM 这些推理引擎存在的全部理由}：算子融合、CUDA Graph、专用的 GEMV 与注意力 kernel（不复制 KV 的 GQA）。

\section{动手练习（按难度递增）}

\begin{enumerate}
  \item 把 \texttt{repeat\_interleave} 去掉：把 Q reshape 成 \texttt{[B, n\_kv, group, T, hd]} 直接和 \texttt{[B, n\_kv, 1, S, hd]} 的 K 做广播矩阵乘，测 Decode 耗时变化。（第 8 篇 §4.3 说明了为什么这很重要）
  \item 用 \texttt{torch.nn.functional.scaled\_dot\_product\_attention} 替换手写的注意力，确认结果不变，比较 Prefill 耗时（它内部会调用 FlashAttention，见第 9 篇）。
  \item 加入第 7 篇的 Temperature + Top-P 采样，并用 \texttt{torch.manual\_seed} 与 HF 的 \texttt{generate(do\_sample=True, ...)} 对比分布是否一致（提示：逐 token 对拍需要完全相同的随机数使用顺序，比较难；改为比较多次采样的统计分布）。
  \item 支持 batch > 1：多个不同长度的请求一起 Decode，需要 padding 与注意力掩码。体会为什么这件事很麻烦——这就是 Lab 12a（连续批处理）和 Lab 12b（PagedAttention）要解决的问题。
  \item （进阶）参考 \href{https://github.com/pytorch-labs/gpt-fast}{gpt-fast} 用 \texttt{torch.compile(mode="reduce-overhead")}（内部使用 CUDA Graph）加速 Decode，看能逼近带宽下界的多少。
  \item （进阶）把权重换成 Lab 11 里的分组 INT4 量化再反量化回 BF16，测量 PPL 与生成质量的变化（第 11 篇）。
\end{enumerate}

\section{学完之后}

你已经能读懂任何 Llama 系模型的 \texttt{modeling\_*.py}，也知道性能差距在哪。下一步：
\begin{itemize}
  \item {}\textbf{第 9 篇：FlashAttention} + \textbf{Lab 09}
  \item {}\textbf{第 12 篇：推理服务系统} + \textbf{Lab 12b}
\end{itemize}



\section{实验核心源码精选与解析}

本实验配套有完整可运行的工业级验证代码，重点实现细节如下：


\subsection{源码实现：\texttt{qwen\_from\_scratch.py}}

\lstinputlisting[language=Python, caption={\texttt{qwen\_from\_scratch.py}}]{code/lab07b_qwen_from_scratch_qwen_from_scratch.py}

